\entry{30 October}{Matching the research video's colours}{Communication revision using the existing development-pose render.}

The researcher requested the same palette as the previous week's
contextual-similarity animation. Its retained source supplied the exact
near-black, warm white, muted grey, sage and cranberry colours. The clean
pose and trace use sage; the three corrupted inputs use cranberry.
The four-panel layout and headline-free crop remain the same.

The local renderer reused the retained arrays without changing geometry,
masking, trace gaps or playback. Source and geometry checks matched the
earlier render, and all 360 frames of the 18-second export decoded.
The cover was visually inspected; the lowest text/background contrast was
4.94:1. The earlier media and journal figure retained their original hashes.
The revised post explains the broader one-shot research question and its
controlled training comparison. No scientific run or social post was made.

\evidence{Final local attachment:
\path{output/real-pose-progress/linkedin-wide-palette.mp4}.
Palette provenance, hashes and checks:
\path{output/real-pose-progress/validation/linkedin-palette-checks.json}.
The media and pose arrays remain outside Git.}
